\Subhead{4.9.1}{Взаимодействие волн пространственного заряда с акустическим полем. Акустоэлектрический эффект. Усилители, линии задержки, фильтры, конвольверы и запоминающие устройства}

В пьезоэлектрике упругая деформация создаёт электрическое поле. Основные линейные связи
\begin{equation}
 \sigma_{ij}=c^E_{ijkl}S_{kl}-e_{kij}E_k,\qquad
 D_i=e_{ikl}S_{kl}+\varepsilon^S_{ij}E_j
\end{equation}
совместно с движением среды и законом Гаусса связывают звук и заряд. Подвижные носители экранируют поле волны с задержкой: $\bm j\simeq\sigma\bm E+\rho_q\bm v_d-D\nabla\rho_q$ вместе с непрерывностью заряда определяют фазу этого отклика. Здесь $\rho_q$ — возмущение плотности заряда, $\bm v_d$ — дрейф, $D$ — коэффициент диффузии.


\paragraph{Волна пространственного заряда.}
В одномерной модели с деформацией $S=u_x$ и $D_x=\varepsilon E+eS$ линеаризация непрерывности для возмущения заряда $\rho_q$ даёт
\begin{equation}
 [-\ii(\omega-kv_d)+Dk^2]\rho_q+\ii k\sigma E=0,\qquad
 \rho_q=\ii k(\varepsilon E+eS).
\end{equation}
Исключение $\rho_q$ показывает, как задержка экранирования меняет упругий модуль:
\begin{equation}\label{eq:spacechargeacoustic}
 E=-\frac e\varepsilon\frac{Dk^2-\ii\omega'}{Dk^2+\tau_M^{-1}-\ii\omega'}S,
 \quad c_{\rm eff}=c^E+\frac{e^2}{\varepsilon}
 \frac{Dk^2-\ii\omega'}{Dk^2+\tau_M^{-1}-\ii\omega'},
 \quad\omega'=\omega-kv_d.
\end{equation}
$e$ здесь пьезомодуль, не заряд электрона; $\tau_M=\varepsilon/\sigma$ --- время максвелловской релаксации. Уравнение звука $\rho_m\omega^2=k^2c_{\rm eff}$ содержит комплексную поправку. Для оценки слабого энергообмена $k$ в этой поправке берут вещественным из невозмущённой акустической волны. При положительной $\omega'$ она даёт поглощение, при смене знака возможна передача энергии от дрейфа звуку. Вывод предполагает малую деформацию, один тип носителей, постоянные коэффициенты и одномерную геометрию, без магнитного поля.
% Самостоятельный вывод из линейных пьезосвязей, j=sigma E+rho_q v_d-D grad rho_q,
% непрерывности и закона Гаусса. Формула не приписана конкретной схеме курса.

Без дрейфа проводимость обычно поглощает акустическую энергию. Если проекция дрейфовой скорости на направление волны больше её фазовой скорости, носители могут передавать энергию звуку; источник — постоянное электрическое питание. В сопутствующем поле важна частота $\omega-\bm k\cdot\bm v_d$. Реальная граница усиления учитывает решёточные потери, диффузию и экранирование; одного неравенства скоростей недостаточно для положительного полного усиления.


\paragraph{Прямой акустоэлектрический эффект.}
Распространяющийся звук передаёт носителям импульс и создаёт постоянный ток либо, в разомкнутом образце, напряжение. Пусть $\Gamma_e$ --- затухание \emph{интенсивности за счёт носителей}, $I_{\rm ac}$ --- интенсивность звука. В простой слабодисперсионной модели сила на единицу объёма $F_e=\Gamma_eI_{\rm ac}/v_s$. Баланс импульса $nm_*v/\tau_m=F_e$ даёт
\begin{equation}\label{eq:directAE}
 j_{\rm AE}=\frac{q\tau_m}{m_*}\frac{\Gamma_eI_{\rm ac}}{v_s},\qquad
 \sigma E_{\rm dc}+j_{\rm AE}=0\quad\text{при разомкнутой цепи}.
\end{equation}
Здесь $q$ --- заряд носителя ($-e_0$ для электронов), $\tau_m$ --- время релаксации импульса. Нельзя подставлять в $\Gamma_e$ полные решёточные потери. Физика передачи импульса подтверждена работой Esslinger и соавт.~\cite{esslinger}; приведённый одномерный баланс --- самостоятельное упрощение, без квантового эффекта Холла.

\paragraph{Линии задержки и фильтры.}
Поверхностная акустическая волна удерживается у границы на глубине порядка длины волны. Встречно-штыревой преобразователь создаёт периодическое электрическое поле; акустический синхронизм связывает период рисунка с длиной волны. Два преобразователя на расстоянии $L$ образуют линию задержки $\tau=L/v_s$. Малость $v_s$ относительно $c$ даёт большие задержки в компактном устройстве.

Апертуры, число и веса электродов формируют импульсную характеристику: принимаемый сигнал — взвешенная сумма задержанных вкладов. Это физическая реализация фильтра, включая согласованный. Дисперсионная линия даёт частотно-зависимую задержку и сжатие чирпа. 
\paragraph{Конвольвер и память.}
Две встречные огибающие $f(t-x/v_s)$ и $g(t+x/v_s)$ перемножаются нелинейным взаимодействием, интегрирующий электрод длины $L$ суммирует вклад:
\begin{equation}
 \begin{split}
 V_{\rm out}(t)&\propto\int_{-L/2}^{L/2}f(t-x/v_s)g(t+x/v_s)\dd x\\
 &=v_s\int_{t-L/(2v_s)}^{t+L/(2v_s)}f(u)g(2t-u)\dd u.
 \end{split}
\end{equation}
Если окно охватывает оба сигнала, это свёртка с временем $2t$. Для корреляции нужна соответствующая обращённая/сопряжённая опора. Линейная задержка без перемножения не является конвольвером.

Обычная линия хранит сигнал лишь во время пробега. В варианте акустоэлектронной памяти сигнал записывает пространственное распределение заряда в диодной структуре; оно сохраняется после прохода волны и считывается следующей ПАВ. Существуют конструкции с массивами диодов Шоттки или p--n-переходов, а также с захватом заряда в слое ZnO~\cite{weng}. Это разные реализации, не универсальная схема любой памяти. Проверена первичная университетская аннотация Weng; полный текст и точная кафедральная конструкция не установлены.
